% Authored research/source audit; not an included theorem or a main-gate claim.
% Created 2026-11-20 after checkpoint22; primary originals remain outside archive.
\section*{The equivariant normal-projection step: exact source scope}

\textbf{Status.}\quad
The rank-four equivariant projection/desuspension step remains unproved.
The sources below do not establish zero Casson--Sullivan sufficiency for
the PL image in \(\pi_0P_{\mathrm{TOP}}(\mathbb{RP}^3)\), nor do they
compute the PL-to-TOP component kernel.  This is a source audit and a
statement of the unsupplied implication, not a numbered result or a
claim that the implication is false.

\textbf{The precise interface.}\quad
An odd rank-four joint-germ clutching family over \(S^2\) supplies a
\(C_2\)-microbundle \(\xi\), with trivial action on the base and sign
action in the fibre charts.  Its total space has dimension six and its
fixed zero section has dimension two.  Even if an equivariant smoothing
of the total pair has been constructed, a smooth tubular neighbourhood
of that fixed section gives a potentially different projection.
The desired conclusion is an equivariant microbundle isomorphism, on
one neighbourhood of the complete zero section,
\[
 H:E(\xi)\supset U\longrightarrow E(\nu),\qquad
 \pi_\nu H=\pi_\xi,\qquad H|_{S^2}=\operatorname{id},
 \qquad H(-v)=-H(v),
\]
with the required relative clutching data.  A homeomorphism of the total
pairs alone supplies neither the displayed equality of projections nor
the relative family statement.  The ordinary rank-four low-type
calculation already recorded in the project does not establish this
equivariant assertion.

\textbf{Hirsch: stable, rather than rank-preserving, uniqueness.}\quad
M.~W.~Hirsch, \emph{On normal microbundles}, Topology 5 (1966),
229--240, original at
\url{https://www.maths.gla.ac.uk/~mpowell/Hirsch-normal-microbundles.pdf}.
The definitions on pp.~230--231 include the projection in the normal
microbundle and require the terminal isomorphism to commute with it.
Theorem~3(c), p.~232, supplies uniqueness over \(S^m\) after adding
\(\varepsilon^{m+2}\), relative to the stated common germ.  Thus its
general sphere assertion adds four trivial directions for \(S^2\).
It is neither an unstabilized rank-four conclusion nor an equivariant
theorem.  The present primary full-text opens were
\texttt{turn934028886view1} and \texttt{turn278649275view2}.

\textbf{Rourke--Sanderson: lifts, stable ranges, and a printed discrepancy.}\quad
C.~P.~Rourke and B.~J.~Sanderson, \emph{On topological neighbourhoods},
Compositio Mathematica 22 (1970), 387--424, original at
\url{https://www.numdam.org/item/CM_1970__22_4_387_0.pdf}.
Theorem~5.1, p.~410, classifies normal bundle structures by lifts to
the indicated open-, micro-, or disc-bundle classifying spaces; it
does not make such a lift automatic.  Corollary~5.4, p.~411, has
\(q\ge6\) for the disc comparison and \(q\ge5\) for the open comparison.
The micro/open clause of Corollary~5.5 requires \(q\ge5\).
There is a literal source discrepancy: the disc clause of Corollary~5.5
prints \(q\ge3\), \(n<q\), with uniqueness for \(n<q-1\), whereas the
introduction, p.~388, says codimension six is needed for disc bundles.
The printed ``3'' was checked in an actually rendered and viewed image;
it is not merely a text-extraction error.  No correction or rank-four
consequence is assumed here.  In any case the statements are ordinary,
not \(C_2\)-equivariant.
Original text: \texttt{turn784062641view0}, earlier
\texttt{turn934028886view0}; theorem locators were also inspected through
\texttt{turn357168925view0--3}.

The temporary original PDF is
\texttt{/tmp/Rourke\_Sanderson\_1970\_neighbourhoods.pdf}, SHA256
\texttt{bae3ed103e3a64fe0835908f411713bd7ecf1aad64bed2b0630a07b47b8103db}.
The actually viewed p.~411 image is
\texttt{/tmp/RS1970\_printed411.png}, SHA256
\texttt{bac9cbbb06f1562e21821c0bc8dff13b82bded108fbe26cfc509bcb467fa1997}.
The web screenshot result \texttt{turn259034523view0} itself returned
only a placeholder; the visual check used the local image instead.

\textbf{Stern: the parameter projection is part of the hypothesis.}\quad
R.~J.~Stern, \emph{Classification theorems for parameterized families
of smooth or piecewise linear manifold structures respecting a
submanifold}, Illinois Journal of Mathematics 21 (1977), 190--219,
original author PDF at
\url{https://bpb-us-e2.wpmucdn.com/faculty.sites.uci.edu/dist/3/246/files/2011/03/16_Parameterized.pdf}.
The definitions on pp.~191--192 make a sliced structure a relative CAT
submersion over the specified parameter space.  Properness in
Remark~1.3 does not first turn a given TOP projection into a PL or smooth
submersion.  Theorems~1.6--1.7, p.~196, preserve that parameter
projection, but require the ambient fibre manifold \(M\) to have
dimension at least five, with their additional boundary conditions.
Putting \(M=\mathbb R^4\) therefore fails this hypothesis.  Instead
using the six-dimensional total space changes what is counted as the
fibre and does not preserve the original \(S^2\)-projection.
The embedding applications in \S5 likewise do not add an equivariance
condition.  Full primary text: \texttt{turn784062641view1}; exact
Theorems~1.6--1.7: \texttt{turn187058286view3}.

\textbf{Le Dimet: equivariant globalization keeps the given projection.}\quad
J.-Y.~Le Dimet, \emph{G-vari\'et\'es topologiques et G-microfibr\'es},
Cahiers de topologie et g\'eom\'etrie diff\'erentielle 17 (1976),
397--435, original at
\url{https://www.numdam.org/article/CTGDC_1976__17_4_397_0.pdf}.
The paper treats compact Lie groups and locally linear topological
actions.  Its \S4, Theorem~5, p.~420, is the homotopy theorem for
pullbacks of a supplied G-microbundle.  The equivariant Kister--Mazur
theorem is \S5, Theorem~16, pp.~428--430, for a locally compact,
second-countable base.  Crucially, \S5, Definition~4, pp.~421--422,
defines containment by an equivariant embedding compatible with
\emph{both} the existing projection and the section.  Consequently
the theorem globalizes one microbundle germ; it does not identify two
normal projections on the same total neighbourhood.  It does not
remove the interface stated above.  Primary original opens:
\texttt{turn278649275view0}, \texttt{turn187058286view0--1}.

\textbf{An incomplete isovariant draft is not a replacement.}\quad
The public file
\url{https://math.ucr.edu/~res/inprogress/eqiso01.pdf},
\emph{Equivariant normal invariants and isovariant homotopy
equivalences}, contains four introductory pages and ends with an
unfinished first section.  It was opened as
\texttt{turn961334980view0}; no complete proof was obtained from it.
Its announced uniqueness hypothesis on a product uses
\(\dim F+2\le(\dim M+1)/2\), which fails for \((\dim M,\dim F)=(6,2)\).
More fundamentally, isovariant homotopy equivalence alone is not the
projection-preserving microbundle isomorphism required here.  This
draft is recorded only to prevent an unsupported invocation.

The complete published comparison is R.~Schultz,
\emph{Isovariant mappings of degree 1 and the Gap Hypothesis},
Algebraic \& Geometric Topology 6 (2006), 739--762,
\url{https://msp.org/agt/2006/6-2/agt-v6-n2-p11-p.pdf},
opened as \texttt{turn623887482view0}.  Theorem~2.1(i), p.~743,
and its proof, p.~747, yield an equivariant \emph{fibre homotopy}
equivalence of the normal sphere bundles.  Theorem~2.2, p.~743,
deforms equivariant homotopy equivalences to isovariant homotopy
equivalences under the stated smooth, closed, oriented and Gap
Hypothesis assumptions.  Neither conclusion is a fibrewise
homeomorphism.  A total-space equivariant homeomorphism is already
isovariant, so merely obtaining isovariance would not improve the
projection interface.  The dimension pair \((6,2)\) meets the
ordinary Gap Hypothesis on p.~739, but its product with an interval
does not.  The publication discusses the stronger product condition
for uniqueness on pp.~744--745; no rank-four projection theorem is
inferred from it.

\textbf{Conclusion of the source check.}\quad
None of these inspected originals supplies the equivariant,
unstabilized, relative rank-four projection identification.  Nor does
failure of these particular hypotheses disprove such a statement.
The actual PL-image and component-kernel questions for the supplied
\(\mathbb{RP}^3\)-concordance remain unchanged.  No original row,
review credit, or included theorem is modified by this note.
